\chapter{Data and Code Availability}
\label{app:availability}

All source code, configuration files, Slurm job scripts, and analysis notebooks
used to produce the results reported in this thesis are publicly available at:

\begin{quote}
\url{https://github.com/hasnat23/master-thesis-rlaif-gca}
\end{quote}

The repository is organized as follows. \url{src/} contains the pipeline
implementation, split into the modules referenced throughout
Chapter~\ref{ch:implementation}: \url{src/data/} (subset selection and record
schemas), \url{src/generation/} (candidate summary generation),
\url{src/judging/} (the AlignScore judge wrapper and both preference-construction
regimes), \url{src/reward_model/} (Bradley--Terry training and cross-validation),
and \url{src/eval/} and \url{src/analysis/} (evaluation and statistics).
\url{configs/} holds the YAML configuration files for subset selection,
generation, and judging. \url{slurm/} holds the batch scripts submitted on MOGON
NHR, including \url{slurm/train_rm_1000.sh}, which produced the headline
$n=1{,}000$ result. \url{analysis/} holds the standalone scripts used for the
formula sweep, the judge-mode sweep, and the holistic-versus-GCA comparison,
together with \url{analysis/aggregate_campaigns.py}, which computes the
run-level means, bootstrap intervals and Wilcoxon tests for each seed campaign,
and \url{analysis/effect_sizes_and_equivalence.py}, which computes the effect
sizes, equivalence tests and minimum detectable effects reported in
Chapter~\ref{ch:results}, Section~\ref{sec:results-effect-sizes}. Both read the
per-run JSON summaries under \url{outputs/} directly, so every reported figure
can be regenerated from the stored run artifacts without retraining.
\url{progress-updates/} contains the dated working notes kept over the course of
the project.

The underlying corpus, CNN/DailyMail version \texttt{3.0.0}
\citep{hermann2015teaching, nallapati2016abstractive}, is publicly available
through the Hugging Face \texttt{datasets} library and is not redistributed in the
repository. The candidate summaries, preference files, and trained reward-model
checkpoints generated during this project are likewise not committed, as they run
to several gigabytes; they can be regenerated from the corpus using the
configuration recorded in Appendix~\ref{app:reproduction}. The factuality
evaluator, \texttt{yzha/AlignScore} \citep{zha2023alignscore}, and the candidate
generator, Mistral-7B-Instruct-v0.3, are both publicly released model checkpoints
obtained from the Hugging Face Hub.

\chapter{Reproduction Configuration}
\label{app:reproduction}

This appendix records the exact configuration required to reproduce the headline
$n=1{,}000$ result. It restates in executable form what
Chapter~\ref{ch:methodology}, Table~\ref{tab:methodology-summary} summarizes and
Chapter~\ref{ch:implementation} describes in prose. Under the corrected,
fully-seeded training procedure described in
Chapter~\ref{ch:implementation}, Section~\ref{sec:impl-rm-training}, a single
run of these commands at a given seed value is deterministic on fixed
hardware; running the full twenty-run headline campaign means repeating the
reward-model training command below once per seed in $\{1, \ldots, 20\}$, as
\url{slurm/seed_campaign.sh} does.

\section{Summarization Prompt}
\label{app:repro-prompt}

Both candidate summaries for a given article are generated from the same fixed
prompt template (\texttt{SUMMARIZATION\_PROMPT} in
\url{src/generation/candidates.py}); the two candidates differ only in the
sampling parameters applied at decoding time.

\begin{quote}
\small\ttfamily
Summarize the following news article in a concise paragraph.\\[2mm]
Article:\\
\{article\}\\[2mm]
Summary:
\end{quote}

\section{Pipeline Parameters}
\label{app:repro-params}

Table~\ref{tab:app-params} lists the parameters of each pipeline stage as used for
the final comparison. Where a value differs from the default recorded in the
corresponding YAML file, the command-line override in
Section~\ref{app:repro-commands} is authoritative; this is the case for the
subset size, the tie margin, and the GCA exponent $\alpha$, each of which changed
over the course of the project as described in Chapter~\ref{ch:methodology},
Section~\ref{sec:methodology-evolution}.

\begin{table}[htbp]
\centering
\caption[Pipeline parameters for the final $n=1{,}000$ comparison]{Pipeline
parameters for the final $n=1{,}000$ comparison.}
\label{tab:app-params}
\begin{tabularx}{\textwidth}{llX}
\toprule
Stage & Parameter & Value \\
\midrule
\multirow{5}{*}{Subset selection}
  & Corpus              & CNN/DailyMail 3.0.0, \texttt{test} split \\
  & Samples             & 1{,}000 \\
  & Subset seed         & 200 \\
  & Max article length  & 8{,}000 characters \\
  & Max summary length  & 2{,}000 characters \\
\midrule
\multirow{6}{*}{Candidate generation}
  & Model               & Mistral-7B-Instruct-v0.3, \texttt{bfloat16} \\
  & Candidate A         & $T=0.7$, top-$p=0.9$ \\
  & Candidate B         & $T=1.0$, top-$p=0.95$ \\
  & Max input tokens    & 1{,}024 \\
  & Max new tokens      & 256 \\
  & Batch size          & 4 \\
\midrule
\multirow{4}{*}{Preference construction}
  & Evaluator           & \texttt{yzha/AlignScore}, \texttt{FacebookAI/roberta-base} \\
  & Evaluation mode     & \texttt{nli} \\
  & GCA exponent        & $\alpha = 0.0$ (simple mean) \\
  & Tie margin          & $0$ \\
\midrule
\multirow{7}{*}{Reward-model training}
  & Backbone            & \texttt{FacebookAI/roberta-base} \\
  & Epochs              & 5 \\
  & Batch size          & 8 \\
  & Learning rate       & $2\times10^{-5}$, linear warmup over 10\% of steps \\
  & Max sequence length & 512 tokens \\
  & Max article chars   & 2{,}000 \\
  & Cross-validation    & 5-fold; training seeds 1--20 (headline campaign,
    Chapter~\ref{ch:results}, Table~\ref{tab:seed-campaign-full}) \\
\bottomrule
\end{tabularx}
\end{table}

\section{Command-Line Invocations}
\label{app:repro-commands}

The three stages are run in sequence. Preference construction produces both
conditions from a single invocation, and reward-model training consumes both
preference files so that the holistic and GCA models are trained under an
identical recipe within one process.

\begin{quote}
\footnotesize
\begin{verbatim}
# 1. Subset selection
python scripts/01_prepare_subset.py --n-samples 1000 --seed 200

# 2. Candidate generation
python scripts/02_generate_candidates.py \
    --config configs/generation_1000.yaml

# 3. Preference construction (both conditions)
python src/judging/build_reward_preferences.py \
    --mode both --alpha 0.0 --margin 0 \
    --alignscore-evaluation-mode nli \
    --alignscore-backbone FacebookAI/roberta-base

# 4. Reward-model training (5-fold CV, both conditions)
python src/reward_model/run_training.py \
    --holistic data/preferences_1000/\
holistic_reward_preferences_1000.jsonl \
    --gca      data/preferences_1000/\
gca_reward_preferences_1000.jsonl \
    --output-dir outputs/reward_models_1000 \
    --backbone FacebookAI/roberta-base \
    --epochs 5 --batch-size 8 --lr 2e-5 \
    --max-length 512 --max-article-chars 2000 \
    --kfold 5 --seed 1
\end{verbatim}
\end{quote}

\noindent
This reproduces one run of the headline campaign, seed 1 of 20. Substituting
each of \texttt{--seed 2} through \texttt{--seed 20} in turn reproduces the
remaining nineteen.

\noindent
In step 4 the two preference-file paths are shown split across lines for
typesetting only; the trailing backslash inside a path is a line continuation in
this listing, not part of the filename.

On the MOGON NHR cluster these are submitted as Slurm batch jobs
(\url{slurm/generate_candidates_1000.sh}, \url{slurm/build_preferences_1000.sh},
\url{slurm/train_rm_1000.sh}) on partition \texttt{a100dl} with one
A100-SXM4-40GB GPU, 32\,GB of host memory, and four CPU cores per task.

\chapter{Example Preference Record}
\label{app:example-record}

Every record written by the GCA preference-construction stage retains the full
sentence-level detail behind the summary-level decision, which is what makes the
qualitative disagreement analysis in Chapter~\ref{ch:results} possible. The
listing below shows one such record, abridged: the article, reference summary, and the
\texttt{chosen}/\texttt{rejected} fields are truncated, and the sentence list for
summary B is cut after the first three sentences. All numeric values are
reproduced verbatim.

The record shown was produced by the earlier $n=200$ run, which used the
exploratory settings $\alpha = 0.5$ and margin $=0.05$ rather than the final
$\alpha = 0.0$ and margin $=0$ (Chapter~\ref{ch:methodology},
Section~\ref{sec:methodology-evolution}). It is included to illustrate the record
\emph{schema} and the relationship between the sentence scores and the aggregate,
not as an instance of the final configuration. Note in particular that
\texttt{gca\_score\_a} $=0.382$ is well below \texttt{mean\_score\_a} $=0.629$:
this is the penalty term at $\alpha = 0.5$ pulling the aggregate down towards the
summary's weakest sentence, the behaviour that Chapter~\ref{ch:results},
Section~\ref{sec:results-formula-opt} found to perform worse than the simple mean.

\begin{quote}
\footnotesize
\begin{verbatim}
{
  "sample_id": "cnn_00000_b0b34b2d843e",
  "article": "LONDON, England (Reuters) -- Harry Potter star Daniel ...",
  "reference_summary": "Harry Potter star Daniel Radcliffe gets ...",
  "summary_a": "Daniel Radcliffe, the 18-year-old actor who plays ...",
  "summary_b": "On Monday, July 23, 2007, Harry Potter actor Daniel ...",
  "gca_score_a": 0.382102,
  "gca_score_b": 0.169184,
  "gca_score_diff": 0.212918,
  "mean_score_a": 0.629039,
  "mean_score_b": 0.578186,
  "min_score_a": 0.232103,
  "min_score_b": 0.049505,
  "n_sentences_a": 3,
  "n_sentences_b": 6,
  "n_low_faithfulness_a": 1,
  "n_low_faithfulness_b": 3,
  "decision": "A",
  "judge_method": "gca",
  "judge_model": "yzha/AlignScore [FacebookAI/roberta-base]",
  "margin": 0.05,
  "alpha": 0.5,
  "sentence_details_a": [
    { "index": 0,
      "text": "Daniel Radcliffe, the 18-year-old actor who plays Harry
               Potter, has gained access to a 20 million pound fortune
               on his 18th birthday.",
      "score": 0.748840, "label": "faithful" },
    { "index": 1,
      "text": "Despite the newfound wealth, he has stated that he does
               not plan to be extravagant, and will instead use his
               money to buy books, CDs, and DVDs.",
      "score": 0.232103, "label": "hallucinated" },
    { "index": 2,
      "text": "Radcliffe has filmed a TV movie and an Australian film,
               and has made his stage debut in \"Equus.\" He is expected
               to face closer media scrutiny now that he is legally an
               adult.",
      "score": 0.906173, "label": "faithful" }
  ],
  "sentence_details_b": [
    { "index": 0, "text": "On Monday, July 23, 2007, ...",
      "score": 0.475803, "label": "hallucinated" },
    { "index": 1, "text": "Radcliffe has stated he has no ...",
      "score": 0.997014, "label": "faithful" },
    { "index": 2, "text": "Instead, he plans to use the money ...",
      "score": 0.049505, "label": "hallucinated" }
  ]
}
\end{verbatim}
\end{quote}

Two features of this record are worth drawing out, because they illustrate
mechanisms discussed in the main text. First, the \texttt{label} field is a
threshold applied at a score of 0.5 for reporting convenience only; it plays no
part in the aggregation, which operates on the continuous scores. Second, summary
B contains both the highest-scoring sentence in the pair (0.997) and the
lowest (0.050). Under the holistic condition this internal variation is not
observable at all: the summary receives a single score. This is precisely the
information that sentence-level decomposition exposes and holistic scoring
discards, and whether exposing it yields a more learnable preference signal is the
question the thesis sets out to answer.
